\section{Five numerical patterns in graph towers}\label{sec:graphs}
DuBose and Vallières prove that the valuation of the number of spanning trees
in a $\mathbb Z_p^d$-tower of finite graphs eventually agrees with a polynomial
in $p^n$ and $n$. Section 7 of their paper fits candidate polynomials to finite
data for five two-dimensional towers \cite{DV}. The existence theorem alone
does not specify those polynomials or the first level where they hold.
The following computations determine both for the five examples.

For a list $S$ of integer vectors, use the group $G_n=(\mathbb Z/p^n)^d$ as
vertices. For each $s\in S$ and $v\in G_n$ put an undirected edge
between $v$ and $v+s$. Retain parallel edges, including both edges if a
step has order two. Assume the reductions of $S$ span $\mathbb F_p^d$,
so these graphs are connected. Write $T_n=v_p(\kappa_n)$, with $\kappa_n$
the number of spanning trees and $v_p(p)=1$.

\begin{theorem}\label{thm:five}The five towers of \cite[\S7]{DV}
have the following valuations. Write $e_1=(1,0)$ and $e_2=(0,1)$.
\begin{center}\small
\begin{tabular}{ccll}
\toprule
Example & $p$ & generators $S$ & $T_n$\\
\midrule
1 & 2 & $e_1,e_2$ & $(2n+4)2^n-6n-1$ \\
2 & 2 & $e_1,e_1,e_2,e_2$ & $4^n+(2n+4)2^n-6n-2$ \\
3 & 2 & $(1,5),(0,3),(1,2),(0,1)$ & $(n+33/4)2^n-4n-1$ \\
4 & 3 & $e_1,e_2$ & $4\cdot3^n-2n-4$ \\
5 & 3 & $(1,0),(2,3),(1,1)$ & $20\cdot3^{n-1}-2n-8$ \\
\bottomrule
\end{tabular}
\end{center}
Rows 1, 2, 4 and 5 hold for every $n\ge1$. Row 3 holds for $n\ge4$,
and its first three values are $5,19,65$.
\end{theorem}

\subsection{Reading a character at its conductor}
Fix a primitive $p^h$-th root $\zeta$ in $\overline{\mathbb Q}_p$, set
$e=\varphi(p^h)$ and $t=\zeta-1$. The polynomial $\Phi_{p^h}(1+T)$
is Eisenstein, so $v_p(t)=1/e$. This is the exact radius calculation used
by the repository's torsion specialization argument. The graph application
uses this identity at \emph{every} character of a given conductor.

For $u\in(\mathbb Z/p^h)^d$ primitive (at least one coordinate is a unit),
the Laplacian eigenvalue at that character is
\[
 F_u(\zeta)=\sum_{s\in S}(2-\zeta^{\langle u,s\rangle}
                                  -\zeta^{-\langle u,s\rangle}).
\]
Fourier characters diagonalize translations. Over the complex numbers all
these eigenvalues are positive for nontrivial characters by connectedness.
The matrix-tree theorem gives the product of nonzero eigenvalues divided by
$p^{dn}$. Embedding the algebraic identity into $\overline{\mathbb Q}_p$
and grouping by conductor yields
\begin{equation}\label{eq:shells}
 T_n=-dn+\sum_{h=1}^n W_h,\qquad
 W_h=\sum_{[u]} e\,v_p(F_u(\zeta)).
\end{equation}
The brackets denote orbits of primitive vectors under multiplication by units
modulo $p^h$. The action is free and its members are cyclotomic conjugates,
all of the same valuation. This orbit decomposition is already used by
DuBose and Vallières. Our evaluations of the individual orders give the
all-level formulas.

Two elementary rules will be useful. If an integral power series in $T$
has first unit coefficient at degree $k$, and $k<e$, then its value at
$t$ has valuation $k/e$. The same conclusion holds for $k<e+2$ if the
constant and linear coefficients vanish identically. Both assertions follow
by comparing the valuations term by term: lower coefficients divisible by
$p$ contribute above the candidate, and higher powers do also. For arbitrary
$k$ one can reduce the polynomial modulo $\Phi_{p^h}(1+T)$, obtaining
$\sum_{j<e}c_jT^j$; the normalized order is exactly
$\min_{c_j\ne0}\{e\,v_p(c_j)+j\}$. These integers lie in distinct residue
classes modulo $e$. This gives the exact small-conductor check used below.

\subsection{The dyadic square torus and its doubled edges}
Write $X=\zeta^a$, $Y=\zeta^b$ and $F=4-X-X^{-1}-Y-Y^{-1}$.
The identity
\[
 F=4-(X+Y)(1+(XY)^{-1})
\]
makes the exceptional characters explicit. For $h\ge3$, put $e=2^{h-1}$.
For a mixed-parity primitive pair, both $a-b$ and $a+b$ are odd.
Each factor in the product therefore has valuation $1/e$, giving normalized
order $2$. There are $2e$ mixed-parity orbits, contributing $4e$.

For the remaining orbits normalize $a=1$, so $b$ runs through the $e$ odd
residues modulo $2^h$. The normalized orders of the factors are those of
$1+\zeta^{1-b}$ and $1+\zeta^{1+b}$. For an exponent $k$ put $r=k+e$
modulo $2^h$. If $r\ne0$, this order is $2^{v_2(r)}$. If $r=0$ the factor
vanishes. For either sign, $1\pm b+e$ runs exactly once through the even
residues. Of a given pair of exponents, one is $2$ modulo $4$, so its order
is $2$. The other, when finite, has order at most $e$.
Thus a finite product has normalized order at most $e+2<2e$, and subtracting
from $4$ preserves that order. At a vanishing factor the order of $F$ is
$2e$ instead.

For either sign, the sum over the nonzero even residues is
$\sum_{i=1}^{h-1}(e/2^i)2^i=(h-1)e$. At each of the two exceptional orbits
the other factor has already contributed $2$ to this sum, so each requires a
correction of $2e-2$. The odd--odd contribution is
$2(h-1)e+2(2e-2)$. Consequently
\[
 W_h=(h+3)2^h-4\quad(h\ge3).
\]
At $h=1$ the eigenvalues are $4,4,8$, giving $W_1=7$.
At $h=2$ the four mixed orbits each contribute $2$, and the two odd--odd
orbits each contribute $4$, giving $W_2=16$. Equation~\eqref{eq:shells}
gives row 1 by summation. Doubling each edge multiplies the number of
spanning trees by $2^{4^n-1}$, since a tree uses $4^n-1$ edges with two
independent parallel choices each. This proves row 2.

\subsection{The four-loop dyadic example with unequal slopes}
In Example~(3) of Section~7 of DuBose and Valli\`eres, the four directions are $(1,5),(0,3),(1,2),(0,1)$. Write $T_n=v_2(\kappa_n)$ for its spanning-tree valuation. Then
\[
 T_n=n2^n+\tfrac{33}{4}2^n-4n-1\qquad(n\ge4),
 \qquad (T_1,T_2,T_3)=(5,19,65).
\]
To see this, write $F=8-\sum_{(u,v)}(X^uY^v+X^{-u}Y^{-v})$.
At conductor $2^s$, take $\zeta$ primitive, $e=2^{s-1}$, and let $W_s$ be the sum of $v_2(F(\zeta^a,\zeta^b))$ over pairs modulo $2^s$ with at least one odd entry. The two vectors $(1,2),(0,1)$ generate the group, so the only zero of $F$ among these characters is the excluded trivial character. The matrix-tree theorem and character diagonalization give $T_s-T_{s-1}=W_s-2$.

The decisive local computation is the following identity over $\mathbb F_2$:
\[
 XY^5\overline F=(Y+1)\big[(Y^2+Y+1)(X^2Y^7+1)+XY^2(Y+1)^5\big].
\]
Substitute $X=(1+z)^a$, $Y=(1+z)^b$, using nonnegative representatives of $a,b$. For $b\ne0$, put $B=2^{v_2(b)}$, $C=2^{v_2(2a+7b)}$. In characteristic two, $(1+z)^j-1$ has order $2^{v_2(j)}$. Since $C$ is a power of two, $C\ne5B$. The resulting order of $\overline F$ is exactly $D=B+\min(C,5B)$. Whenever $D<e$ this gives $ev_2(F(\zeta^a,\zeta^b))=D$: in the integral expansion at $z=\zeta-1$ the lower powers have even coefficients, the coefficient of $z^D$ is odd, and $v_2(\zeta-1)=1/e$. One can use the polynomial $XY^5 F$ to avoid negative exponents throughout.

For $s\ge5$, the sums over the following disjoint classes exhaust the conductor shell.
\begin{center}
\begin{tabular}{lll}
\hline
Class & Normalized order $ev_2(F)$ & Contribution to $W_s$\\
\hline
$b$ odd & $2$ & $4e$\\
$b\equiv2\pmod4$ & $6,10,12$ & $17e/4$\\
$v_2(b)=k$, $2\le k\le s-2$ & $2^k+2$ & $(s-3)e+e-4$ (all $k$)\\
$b=0$ & $e+2$ & $e+2$\\
$b=e$ & $2e$ & $2e$\\
\hline
\end{tabular}
\end{center}
Indeed, for fixed $b\equiv2\pmod4$, as $a$ runs through odd residues, $C$ is $4,8,$ or at least $16$ in proportions $1/2,1/4,1/4$. For $k\ge2$ we have $C=2$ and $e/2^k$ eligible values of $b$ at depth $k$. All three generic rows have $D<e$. For the two exceptional rows, the integer identities $F(X,1)=2(2-X-X^{-1})$ and $F(X,-1)=12$ give the listed orders directly. Summation yields
\[
 W_s=(s+37/4)2^{s-1}-2\quad(s\ge5).
\]
The small-shell totals $W_1,W_2,W_3,W_4$ are $7,16,48,116$. For an exact short verification, use one representative per Galois orbit, $(a,1)$ with arbitrary $a$ and $(1,b)$ with $b$ even. Reduce $F(Z^a,Z^b)$ modulo $Z^e+1$ and write the result as $\sum_{j<e}c_j(Z-1)^j$. The orbit contributes $\min_{c_j\ne0}(ev_2(c_j)+j)$, since its size is $e$ and these orders are distinct modulo $e$. At $s=1,2,3,4$, the multisets of orbit contributions, using superscripts for multiplicities, are respectively
\[
 \{2^{[2]},3\},\quad \{2^{[4]},4^{[2]}\},\quad
 \{2^{[8]},6^{[2]},8,12\},\quad
 \{2^{[16]},6^{[4]},10^{[2]},16,24\}.
\]
Thus $T_4=179$. Summing the shell recurrence proves the formula for all subsequent levels.


\subsection{The two ternary towers}
For the square torus $p=3$, the coefficient of $T^2$ in
$F_u(1+T)$ is $-(a^2+b^2)$ and the lower two coefficients are exactly zero.
Every primitive pair has $a^2+b^2\ne0$ modulo three. Thus each orbit has
order $2$, even at $h=1$ by the strengthened comparison rule. The number
of orbits is $4\cdot3^{h-1}$, giving row 4 of Theorem~\ref{thm:five}.

For row 5 the projected steps are $a,,2a+3b,,a+b$. At $h=1$ the
eigenvalue is three times the number of nonzero projected steps modulo
three. The four orbits have normalized orders $4,4,2,2$, so $W_1=12$.
For $h\ge2$, classify the primitive direction modulo nine. The following
table gives the first nonzero degree $k$ of $F_u(1+T)$ modulo $(3,T^9)$.
\begin{center}\small
\begin{tabular}{c|rrrrrrrrr|ccc}
$u$ & $(1,0)$ & $(1,1)$ & $(1,2)$ & $(1,3)$ & $(1,4)$ & $(1,5)$ &
$(1,6)$ & $(1,7)$ & $(1,8)$ & $(0,1)$ & $(3,1)$ & $(6,1)$\\
\hline
$k$&4&6&2&6&4&2&4&4&2&2&2&2
\end{tabular}
\end{center}
For an explicit verification put $J=\{a,2a+3b,a+b\}$ as a multiset.
For $1\le j\le6$ the coefficient is the reduction modulo three of
$-\sum_{k'\in J}(\binom{[k']_9}{j}+\binom{[-k']_9}{j})$,
where the representatives lie in $[0,8]$. The constant is zero before reduction.
The table persists at every higher lift since
$(1+T)^9\equiv1\pmod{3,T^9}$. Normalizing a unit multiple amounts to a
change of variable with invertible linear term. In the original Laurent
expansion both the constant and the linear terms vanish exactly; thus
$k\le6<e+2$ makes each displayed $k$ the normalized valuation at every
$h\ge2$. Every displayed orbit has $3^{h-2}$ orbit lifts. The sum of the
table is $40$, so $W_h=40\cdot3^{h-2}$. Summing from $W_1=12$ proves row 5.

